\chapter{Proposta de experimento}

Para realizar o experimento será necessário treinar um modelo de rede
neural que seja capaz, ou esteja próximo, de decifrar um CAPTCHA. Para
isso serão efetuadas três etapas básicas e comuns quando se trabalha
com redes neurais. Primeiro será coletado o maior número possível de
imagens de CAPTCHA. Em seguida será gerado um dataset com as
características dessas imagens junto com a classe em que pertence. A
partir daí podemos realizar a configuração e treinamento da rede
neural. E por fim será testada a acurácia do melhor modelo mediante
imagens de teste.

\section{Coleta de imagens}

Como o escopo do trabalho não contempla a automatização da recuperação
de informações de websites públicos, foi disponibilizado um
repositório com as imagens necessárias. Esse repositório possui
188.564 imagens e foi disponibilizado pela empresa Neoway. As imagens
se tratam de um CAPTCHA publicado pelo site do SINTEGRA de Santa
catarina
(\url{http://sistemas3.sef.sc.gov.br/sintegra/consulta_empresa_pesquisa.aspx}).

\begin{figure}[H]
\centering
\includegraphics[scale=1]{imagens/exemplo_captcha}
\caption{Um exemplo do CAPTCHA utilizado pelo sistema de consulta do
  SINTEGRA de Santa Catarina.}
\label{fig:exemplo_captcha}
\end{figure}


\subsection{Fonte pública}

Para demonstrar a ineficiência de certas imagens de CAPTCHA foi
escolhido um software Web. Este software do SINTEGRA, fornece dados
públicos de contribuintes mediante consulta via website. O SINTEGRA é
o Sistema Integrado de Informações sobre Operações Interestaduais com
Mercadorias e Serviçoes. Esta fonte pública possui dados fornecidos
pelos próprios contribuintes na hora do cadastro. Os comerciantes ou
profissionais autônomos fazem seu cadastro para facilitar o comércio
de produtos e prestação de serviços. O cadastro contempla inscrição da
pessoa física ou jurídica, endereço e informações complementares
referentes ao fisco estadual.

\section{Geração do Conjunto de dados}

O conjunto de dados (ou ``dataset'') que alimenta a rede neural é
gerado em tempo de execução do treinamento. Cada imagem é lida de seu
diretório em disco e carregada na memória como uma matriz de valores
de pixel. Ao final deste processo há um vetor em memória com todas
imagens existentes já pré-processadas. Isso é feito para o dataset de
treinamento e de teste. O dataset de treinamento terá a maioria das
imagens, portanto {\bf 180 mil}.

\subsection{Pré-processamento}

A fase de pré-processamento das imagens é mínima e é feita junto com a
geração do conjunto de dados.

\begin{itemize}
\item{\bf Escala de cinza}

Ao gerar um \textit{array} representativo da imagem, apenas é
considerado um valor de escala de cinza da imagem, assim padronizando
os valores de intensidade de pixels entre 0 e 1.

\item{\bf Redimensionamento}

Ao gerar o \textit{array} que representa a imagem, é feito um cálculo
para diminuir a imagem com base em uma escala. Essa escala será
configurada à partir de um valor padrão para a largura e altura das
imagens. 

\end{itemize}

\subsection{Conjunto de dados de teste}

Para o treinamento será necessário um conjunto separado para teste que
não possui nenhuma imagem presente no conjunto de treinamento.
O dataset de testes terá uma amostra bem menor que o conjunto de
treinamento. Portanto terá {\bf 8 mil} imagens.

\section{Treinamento}

Após gerado o conjunto de dados, é possível trabalhar no treinamento
do modelo da rede neural. Para isso será usado o Framework {\bf
  \emph{TensorFlow}}\cite{TensorFlow} destinado à \textit{Deep
  Learning} e um script em Python que fará uso das funções
disponibilizadas pela biblioteca do \textit{TensorFlow}. Assim
realizando o treinamento até atingir um valor aceitável de acerto no
conjunto de teste. O resultado do treinamento será um arquivo binário
representando o modelo que será utilizado para avaliação
posteriormente.

\subsection{Infraestrutura}

Com o intuito de acelerar o processo, foi utilizada uma
máquina com {\bf GPU} para o treinamento. A máquina foi adquirida em
uma \textit{Cloud} privada virtal da AWS. A GPU utilizada se trata de
uma \textit{NVIDIA K80} com 2.496 cores e 12GB de memória de
vídeo. Como processador a máquina possui um \textit{Intel Xeon
  E5-2686v4 (Broadwell)} com 4 cores, e ainda possui 61GB de memória
RAM \cite{GPUinstance}.

\subsection{Bibliotecas utilizadas}

Todo o código foi implementado utilizando a linguagem de programação
{\bf \emph{Python}}, e as seguintes bibliotecas foram utilizadas:

\begin{itemize}

\item {\bf TensorFlow}\cite{TensorFlow}: Um Framework implementado em
  \textit{Python} destinado à \textit{Deep Learning}. Proporciona a criação da
  arquitetura e automatização do processo de treinamento de redes
  neurais com \textit{backpropagation}.

\item {\bf NumPy}\cite{NumPy}: Uma biblioteca em \textit{Python}
  criada para computação científica. Possui um objeto de
  \textit{array} com várias dimensões e várias funções sofisticadas
  para cálculos com algebra linear.

\item {\bf OpenCV}\cite{OpenCV}: Uma biblioteca, implementada em
  C/C++, destinada à computação visual. Utilizada para ler imagens em
  disco e realizar o pré-processamento nas mesmas.

\end{itemize}

\section{Avaliação de acurácia}

Para a avaliação, uma nova amostra de imagens será coletada do mesmo
modo que foram coletadas as imagens para treinamento. Essa amostra
terá uma quantidade maior de imagens do que o conjunto de teste.

Com essa amostra de imagens, será feita a execução do teste do modelo
contra cada uma das imagens, assim armazenando uma informação de erro
ou acerto do modelo. Ao final da execução será contabilizado o número
de acertos e comparado com o número total da amostra de imagens para
avaliação. Resultando assim em uma porcentagem que representa a
acurácia do modelo gerado.

